\chapter{Simulation setup, framework and implementation}

This chapter describes the numerical setup and analysis pipeline used to investigate the persistent central density decline identified in the previous chapters. The focus is on reproducibility: each subsection documents the simulation code configuration, initial-condition choices, diagnostic definitions, and implementation details used to test competing numerical and physical explanations.

The workflow combines (i) controlled collisionless $N$-body runs in SWIFT, (ii) systematic parameter variations of the self-gravity solver, (iii) profile and cumulative-mass diagnostics extracted from snapshots, and (iv) shell-based analyses designed to quantify non-local contributions to the dynamics of the innermost region.

\section{SWIFT simulation setup}
\label{sec:framework_swift_setup}

All simulations in this thesis are dark-matter-only and non-periodic, evolved with SWIFT in self-gravity mode (and external gravity when required by specific tests). In the working setup, runs are launched through
\begin{verbatim}
swift --self-gravity [--external-gravity] params.yml
\end{verbatim}
with outputs consisting of snapshots, global statistics, and timestep logs.

The internal unit system is chosen so that
\begin{equation}
1\ \mathrm{length\ unit} = 1\ \mathrm{kpc}, \qquad
1\ \mathrm{velocity\ unit} = 1\ \mathrm{km\ s^{-1}}, \qquad
1\ \mathrm{mass\ unit} \approx 10^{10}\ M_{\odot},
\end{equation}
which allows straightforward interpretation of profile radii (kpc), velocity dispersions (km/s), and enclosed masses in solar units after post-processing conversion.

\subsection{Baseline gravity and time integration choices}

The parameter-study baseline adopts the self-gravity configuration
\begin{equation}
\eta = 0.002, \qquad \mathrm{MAC}=\mathrm{geometric}, \qquad
\theta_{\mathrm{cr}}=0.7, \qquad \epsilon_{\mathrm{fmm}}=10^{-3},
\end{equation}
with minimum and maximum timesteps typically constrained in the interval
\begin{equation}
\Delta t_{\min}=10^{-6}, \qquad \Delta t_{\max}=10^{-2}
\end{equation}

To probe solver sensitivity, each of the following was varied while keeping the others fixed: timestep factor $\eta$, MAC choice (geometric/adaptive), geometric opening threshold $\theta_{\mathrm{cr}}$, and $\epsilon_{\mathrm{fmm}}$. These controlled tests form the foundation of the numerical robustness analysis presented in the results chapter.

\subsection{Output cadence and bookkeeping}

The production analysis relies on regular snapshot dumps that allow profile reconstruction at multiple times. In parallel, SWIFT writes global conservation and bulk quantities in \texttt{statistics.txt}, and timestep activity in \texttt{timesteps.txt}.

\section{Halo initial conditions and parameter choices}
\label{sec:framework_ic_params}

The halo models are generated as collisionless particle realizations with controlled inner slopes to enable cusp/core comparisons under otherwise matched numerical conditions.

Initial conditions are produced with the pNbody code:
\begin{verbatim}
ic_gen_2_slopes --alpha 1 --beta 3 --Rmin 1e-3 --Rmax 50 --r_cutoff 50 \
				--rho0 0.0027658297046461644 --rs 0.828177997155062 \
				--ptype 1 --mass 61664.060524 -t swift -o halo.hdf5
\end{verbatim}
which corresponds to a two-slope halo with tunable inner structure and an outer truncation radius.

Across the broader thesis campaign, the most relevant controlled parameters are:
\begin{itemize}
\item halo type: cusp and core initial conditions,
\item box-size variants for finite-volume sensitivity checks,
\item gravitational softening variants,
\item self-gravity solver parameters ($\eta$, MAC, $\theta_{\mathrm{cr}}$, $\epsilon_{\mathrm{fmm}}$).
\end{itemize}

The practical philosophy is to keep one baseline run and modify one parameter family at a time, so each subsection in the results can be interpreted as a direct hypothesis test rather than a multi-parameter confounded comparison.

\section{Density, shell, and mass diagnostics}
\label{sec:framework_diagnostics}

This section defines the observables used throughout the thesis.

\subsection{Spherically averaged density profiles}

Density profiles are extracted from snapshots with a radial grid up to a user-defined $r_{\max}$ and $N_r$ bins. In practice, the analysis scripts use both moderate and high radial resolutions, including $N_r=500$ in the inner-region-focused runs.

For shell $i$ with mass $\Delta M_i$ and boundaries $(r_{i,\mathrm{in}},r_{i,\mathrm{out}})$, the estimator is
\begin{equation}
\rho_i = \frac{\Delta M_i}{\frac{4\pi}{3}\left(r_{i,\mathrm{out}}^3-r_{i,\mathrm{in}}^3\right)}.
\end{equation}
Profiles are then compared to the initial state using ratio curves $\rho(r,t)/\rho(r,t_0)$, which makes secular flattening or depletion immediately visible.

\subsection{Softening-scale marker in profile plots}

To separate physically interpretable radii from numerically delicate scales, profile plots include a vertical reference at
\begin{equation}
r_{\mathrm{ref}} = 2.8\,\epsilon,
\end{equation}
with $\epsilon$ the run softening length. This marker is not used as a hard cutoff, but as a visual guide when discussing central profile evolution.

\subsection{Cumulative mass profiles}

Cumulative mass diagnostics are computed via
\begin{equation}
M(<r) = \sum_{r_j < r} m_j,
\end{equation}
including optional decomposition by particle type when needed. The implementation supports multiple centering modes (center of mass, median, manual center, and shrinking-sphere center), with shrinking-sphere used in the dedicated cumulative-mass tests because of its robustness to asymmetries and mild off-centering.

\subsection{Velocity and anisotropy diagnostics}

Using the same radial binning, the analysis computes velocity moments and anisotropy indicators, including
\begin{equation}
\beta(r) = 1 - \frac{\sigma_t^2}{2\sigma_r^2},
\end{equation}
as complementary diagnostics for identifying orbital heating and redistribution mechanisms.

\section{Post-processing pipeline and centering}
\label{sec:framework_pipeline}

The post-processing pipeline is designed to keep the density measurements stable against two common sources of noise: snapshot-to-snapshot variance and small centering offsets. These steps are essential for ensuring that the measured trends reflect the simulation rather than the extraction procedure.

\subsection{Snapshot concatenation}

To reduce stochastic noise in the density-evolution curves, consecutive snapshots are concatenated into larger time blocks before the profile comparison is made. This averaging procedure smooths out short-timescale fluctuations that appear in individual outputs, while preserving the secular evolution of the halo.

The key point is that concatenation changes the visual smoothness of the curves, but not the underlying behaviour: the declining inner-density trend remains present after the time blocks are combined. For that reason, concatenation is used only as a noise-reduction step, not as a way of altering the physical interpretation.

\subsection{Shrinking-sphere centering}

Centering is handled independently from the density estimation. Each snapshot is recentered using a shrinking-sphere algorithm, which iteratively refines the halo center by repeatedly restricting the particle set to smaller radii around the current estimate. This method is robust against mild asymmetries and against the gradual drift that can occur in long collisionless runs.

Using the shrinking-sphere center ensures that the innermost bins of the density profile are measured relative to a dynamically meaningful origin. In practice, this removes one of the most obvious numerical artefacts that could otherwise be mistaken for a central density decline.

\section{Shell-based analysis framework}
\label{sec:framework_shells}

To determine whether the central density decline is driven purely by local inner-halo physics or by non-local coupling from the broader halo structure, we decompose the halo into radial shells and quantify each shell's contribution to the gravitational potential and dynamics of the innermost region.

\subsection{Motivation and design principle}

In a spherically symmetric system, the radial force on a particle at radius $r$ depends only on the enclosed mass $M(<r)$ by the shell theorem. However, the gravitational potential depends on all mass, including material at larger radii. This distinction is crucial: outer shells do not change the local force gradient, but they do deepen the potential well and contribute to the escape-speed scale. The shell-based analysis tests whether the central decline is sensitive to the cumulative influence of outer material on the inner potential structure.

The analysis addresses a specific question: if outer shells contribute substantially to the inner potential depth, then a hybrid model that replaces the outer halo ($r > r_{\mathrm{cut}}$) with an analytic potential while retaining the inner region ($r < r_{\mathrm{cut}}$) as particles should reproduce similar evolutionary trends. Conversely, if the central decline persists even when the outer dynamics is replaced by a smooth analytic model, then the effect is either driven by inner-region physics or by the potential structure rather than by dynamical interactions with individual outer particles.

\subsection{NFW truncated-halo model}

The analysis uses a spherically symmetric Navarro--Frenk--White (NFW) halo model \citep{navarro1996structure}, which is a standard benchmark for dark matter halos:

\begin{equation}
\rho(r) = \frac{\rho_s}{(r/r_s)(1 + r/r_s)^2},
\end{equation}

where $\rho_s$ is the characteristic density and $r_s$ is the scale radius. The enclosed mass within radius $r$ is:

\begin{equation}
M(<r) = 4\pi \rho_s r_s^3 \left[ \ln(1+x) - \frac{x}{1+x} \right], \quad x = \frac{r}{r_s}.
\end{equation}

The halo is truncated at a virial radius $r_{\mathrm{vir}}$. The gravitational potential including the truncation is:

\begin{equation}
\Phi(r) = -\frac{G M(<r)}{r} - 4\pi G \rho_s r_s^2 \left[ \frac{1}{1 + r/r_s} - \frac{1}{1 + r_{\mathrm{vir}}/r_s} \right],
\end{equation}

and the radial acceleration (force per unit mass) is:

\begin{equation}
a(r) = \frac{G M(<r)}{r^2}.
\end{equation}

\subsection{Key diagnostics and definitions}

The analysis is built around four central diagnostic quantities, each addressing a different aspect of halo structure:

\begin{description}

\item[\textbf{Enclosed mass and shell mass fraction}]
The enclosed mass $M(<r)$ tracks how the total mass budget grows with radius. The shell mass fraction is the mass in a radial shell divided by the total halo mass, showing the radial distribution of the mass budget. For an NFW profile, the enclosed mass grows steeply near the center but the shell mass fraction increases toward larger radii, indicating that the outer halo contains a substantial fraction of the total mass.

\item[\textbf{Central potential depth from individual shells}]
For each shell at radius range $[r_i, r_{i+1}]$, we calculate its contribution to the gravitational potential at the center ($r = 0$). This is given exactly by:

\begin{equation}
\Delta\Phi_{\rm shell}(0) = -4\pi G \rho_s r_s^2 \left[ \frac{1}{1 + r_i/r_s} - \frac{1}{1 + r_{i+1}/r_s} \right].
\end{equation}

Shells well outside the center still make sizeable contributions to the potential depth at $r = 0$ (and therefore to the escape-speed scale $v_{\rm esc} \propto \sqrt{-2\Phi}$). The fractional contribution of each shell quantifies this effect.

\item[\textbf{Outer-potential fraction}]
For a given probe radius $r_{\mathrm{probe}}$ and cutoff radius $r_{\mathrm{cut}}$, the outer-potential fraction is defined as:

\begin{equation}
f_{\mathrm{outer}}(r_{\mathrm{probe}}, r_{\mathrm{cut}}) = \frac{\left|\Phi_{>r_{\mathrm{cut}}}(r_{\mathrm{probe}})\right|}{\left|\Phi_{\mathrm{total}}(r_{\mathrm{probe}})\right|},
\end{equation}

where $\Phi_{>r_{\mathrm{cut}}}$ is the potential contribution from material outside the cutoff and $\Phi_{\mathrm{total}}$ is the total potential at that radius. This measures how strongly the outer halo influences the local potential at the probe radius. By varying $r_{\mathrm{cut}}$, we can determine the radius beyond which the outer halo becomes a small correction.

\item[\textbf{Force contribution by shell}]
The radial force (acceleration) at a probe radius $r_{\mathrm{probe}}$ is:

\begin{equation}
a(r_{\mathrm{probe}}) = \frac{G M(<r_{\mathrm{probe}})}{r_{\mathrm{probe}}^2}.
\end{equation}

Under spherical symmetry, only mass interior to $r_{\mathrm{probe}}$ contributes. This diagnostic confirms that outer shells, despite deepening the potential, do not directly modify the local force field---a consistency check on the shell-theorem interpretation.

\end{description}

\subsection{Implementation: analytical NFW shell decomposition}

Rather than relying on a numerical computation of the integrals, the shell analysis uses exact analytic formulas for a truncated NFW profile. The halo is decomposed into logarithmically spaced spherical shells:

\begin{equation}
r_i = r_{\mathrm{min}} \left(\frac{r_{\mathrm{vir}}}{r_{\mathrm{min}}}\right)^{i / N_{\rm shells}}, \quad i = 0, 1, \ldots, N_{\rm shells},
\end{equation}

where typically $N_{\rm shells} = 80$ and $r_{\mathrm{min}} \approx 0.02$ kpc.

For each shell $[r_i, r_{i+1}]$, exact closed-form expressions compute: (1) shell mass $M_{\rm shell} = M(<r_{i+1}) - M(<r_i)$, (2) shell contribution to potential at any radius, and (3) cumulative potential contributions for integration over shells.

The analytical approach avoids numerical-integration noise and ensures all diagnostics are computed to high precision. The Python implementation uses these formulas to generate diagnostic plots showing: (i) enclosed mass and shell mass-fraction distribution, (ii) shell contributions to the central potential depth, (iii) outer-potential fraction as a function of cutoff radius, and (iv) a 2D map of outer-fraction versus both probe and cutoff radii.


\section{Numerical implementation of the hybrid model}
\label{sec:framework_hybrid}

Motivated by the shell results, a dedicated hybrid gravity model was implemented and tested on a custom SWIFT branch (\texttt{truncated-NFW}). The model enforces a hard radial split at a cutoff radius $R_{\mathrm{cut}}$:
\begin{itemize}
\item inner region ($r \le R_{\mathrm{cut}}$): standard self-gravity, no external NFW potential,
\item outer region ($r > R_{\mathrm{cut}}$): external NFW potential applied, with optional masking of dark-matter self-gravity sourcing outside the cutoff.
\end{itemize}

The design goal is not to claim a final physical model, but to construct a targeted numerical experiment: if outer-halo coupling is a major driver of central evolution, changing how the outer region sources gravity should produce a measurable response in inner-density trends.

\subsection{Implementation objectives and design choices}

The truncated-NFW backend is designed to isolate the dynamical role of the outer halo while keeping the inner region fully self-consistent. The implementation follows three guiding principles:

\begin{itemize}
\item \textbf{Clear separation of responsibilities}: inner particles (\(r\le R_{\mathrm{cut}}\)) evolve under the usual SWIFT self-gravity solver and provide the only resolved self-gravity sources inside the cutoff. Outer particles (\(r>R_{\mathrm{cut}}\)) evolve in the analytic truncated-NFW field, with resolved gravity contributions removed by construction.
\item \textbf{Minimal invasive changes}: the feature is implemented as an optional, compile-time selectable external-potential backend and does not alter the default gravity code paths when disabled.
\item \textbf{Deterministic, documented boundary convention}: the hard boundary is defined as inner for \(r\le R_{\mathrm{cut}}\) and outer for \(r>R_{\mathrm{cut}}\), avoiding ambiguous behaviour at the cutoff.
\end{itemize}

\subsection{Build-time wiring and backend selection}

The backend is registered as an external-potential option and selected at configure time. The recommended configure invocation is:

\begin{verbatim}
./configure --with-ext-potential=truncated-nfw
\end{verbatim}

At build-time the configure script defines the macro identifying the truncated-NFW backend and compiles the corresponding header implementation. This approach keeps the extra code path isolated behind a single configure flag and reduces the risk of accidental behavior changes in other backends.

Files involved at build time include the project's \texttt{configure.ac} (backend option and macro definition), \texttt{src/potential.h} (backend include branch), and \texttt{src/potential/truncated\_nfw/potential.h} (backend implementation). Gravity-side hooks live in \texttt{src/gravity/*}, \texttt{src/multipole.h}, and mesh-gravity files, but remain dormant unless the truncated-NFW backend is selected.

\subsection{Runtime parameters and YAML configuration}

Runtime control is provided through the existing \texttt{NFWPotential} block in the parameter YAML. The truncated-NFW mode extends the canonical NFW keys with a small set of required and optional keys:

\begin{itemize}
\item \texttt{cutoff\_radius:} required floating-point value; the radial cutoff $R_{\mathrm{cut}}$ in kpc.
\item \texttt{dm\_mask\_outside\_cutoff:} optional boolean (default 1). If true, DM outside the cutoff is excluded as a source of self-gravity.
\item \texttt{position:} center of the external potential (array); default behaviour follows the usual \texttt{NFWPotential:useabspos} semantics.
\item \texttt{useabspos:} boolean as in the standard NFWPotential, controls interpretation of \texttt{position}.
\end{itemize}

The parameter parser in the backend validates that \texttt{cutoff\_radius} is present when truncated-NFW is selected and prints the effective setup at startup (center, halo parameters, cutoff, masking flag). The same center definition is propagated to the masking logic through dedicated globals (mask enabled flag, mask center, and $R_{\mathrm{cut}}^2$), so potential and masking use exactly the same geometric reference.

\subsection{Two-layer masking logic: sources and targets}

The gravity split is enforced through two complementary mechanisms.

\textbf{Source-side masking.} Dark-matter particles outside the cutoff use an effective gravitational source mass of zero through \texttt{gravity\_get\_effective\_mass\_from\_position()}. Operationally this corresponds to \texttt{grav\_mass\_factor}=0 for outer DM and \texttt{grav\_mass\_factor}=1 otherwise. The effective mass is then consumed consistently by pairwise interactions, gravity caches, multipole construction, and no-cache pair loops (including \texttt{runner\_doiact\_grav.c}), so outer DM does not source resolved self-gravity.

\textbf{Target-side suppression.} Outer particles are also prevented from receiving resolved self-gravity updates. This ensures that the outer domain is analytic-only not only as a source prescription but also at the level of target acceleration/potential write-back.

To implement this, a common geometric helper
\texttt{gravity\_is\_outside\_truncated\_nfw\_cutoff(const double x[3])}
was added in both gravity variants:
\texttt{src/gravity/Default/gravity.h} and
\texttt{src/gravity/MultiSoftening/gravity.h}.

This helper is then used to gate resolved-gravity write-back paths:

\begin{itemize}
\item \texttt{src/gravity\_cache.h}: in \texttt{gravity\_cache\_write\_back(...)} a guard skips adding cached pair/tree acceleration and potential to outer targets;
\item \texttt{src/multipole.h}: in \texttt{gravity\_L2P(...)} an early return prevents multipole target updates for outer particles;
\item \texttt{src/mesh\_gravity.c}: in \texttt{mesh\_to\_gpart\_CIC(...)} an early return suppresses local PM interpolation for outer particles;
\item \texttt{src/mesh\_gravity\_mpi.c}: in \texttt{mesh\_patch\_to\_gparts\_CIC(...)} an early return suppresses MPI PM interpolation for outer particles.
\end{itemize}

Operationally, for any target with $r>R_{\mathrm{cut}}$, resolved pair/tree/mesh contributions are not written to the particle state. Outer particles therefore evolve only under analytic external potential terms.

The same predicate is used in both local and MPI-distributed mesh paths, avoiding rank-dependent behaviour at the cutoff surface.

\subsection{External potential behaviour and timestep coupling}

The external-potential API is responsible for returning three quantities for each particle position: acceleration \(\mathbf{a}_{\rm ext}(\mathbf{x})\), potential energy contribution \(\Phi_{\rm ext}(\mathbf{x})\), and optional timestep constraint multipliers. In truncated-NFW mode these are implemented as

\begin{itemize}
\item for $r \le R_{\mathrm{cut}}$: \(\mathbf{a}_{\rm ext}=0\), \(\Phi_{\rm ext}=0\) (no external field acting on inner particles),
\item for $r > R_{\mathrm{cut}}$: standard truncated NFW analytic expressions for \(\mathbf{a}_{\rm ext}\) and \(\Phi_{\rm ext}\) are returned.
\end{itemize}

The analytic potential/acceleration are injected through the standard external-potential hooks (including \texttt{gravity\_add\_comoving\_potential(...)}), while resolved-gravity suppression is handled in gravity and mesh write-back paths.

This separation clarifies potential bookkeeping: outer particles no longer accumulate resolved self-potential from tree/pair/mesh, but they still receive the analytic external contribution via the external backend and enter \(E_{\rm pot,ext}\) diagnostics through \texttt{external\_gravity\_get\_potential\_energy(...)}.

The external potential may optionally supply a timestep multiplier for outer particles (for example to limit timesteps when the analytic field varies strongly), but by default the domain's existing timestep criteria govern integration.

\subsection{Centering, cutoff convention and edge cases}

The backend uses a single, explicitly provided center for both the NFW potential and the cutoff test. Two configuration choices affect this center: absolute coordinates (\texttt{useabspos=1}) or coordinates relative to box origin/periodic convention (\texttt{useabspos=0}). The policy is documented at startup and logged.

We adopt the hard-step convention: particles at exactly $r=R_{\mathrm{cut}}$ are classified as inner (\(r\le R_{\mathrm{cut}}\)). This avoids oscillatory reclassification for particles grazing the boundary and is consistent with the mathematical definitions used in the analytic diagnostics.

\subsection{Serialization, restart safety and diagnostics}

All runtime parameters are recorded in the standard external-potential state so that restarts reproduce the mask and potential choices. On startup the backend prints a compact diagnostics summary that includes backend name, chosen center, $r_{\mathrm{cut}}$, and whether DM masking is active. The initialization path that sets masking state and effective masses is executed at fresh start and restart, and is refreshed at the gravity update cadence (per step or per tree/domain rebuild, depending on run configuration).

Because the modification changes the effective source mass (and therefore total self-gravity source) the global energy bookkeeping routines are instrumented to report the expected non-conservative terms and to show the separation between particle-sourced and analytic-potential contributions to the total potential energy. These diagnostics are intended for validation and are not enabled in tight production runs by default.

\subsection{Foreign-particle and MPI consistency}

To preserve exact force symmetry across MPI boundaries, the effective mass semantics must be identical for both local and foreign particles. The implementation therefore applies one of two consistent strategies when packing remote particles:

\begin{enumerate}
\item export the effective mass (\(m_{\rm eff}\)) directly in the packed structure, or
\item export both the raw mass and the mask-factor, and reconstruct \(m_{\rm eff}\) on the receiving side.
\end{enumerate}

Either strategy yields the same physics if used consistently; the implementation enforces a single semantics for local and remote particles so multipole and pairwise forces remain symmetric across rank boundaries.

\subsection{Validation strategy}

The implementation is accompanied by a small validation suite that exercises the key correctness properties:

\begin{itemize}
\item single-particle checks: place test particles just inside and just outside $R_{\mathrm{cut}}$ and verify that inner particles experience no external acceleration while outer particles feel the analytic NFW acceleration;
\item self-gravity exclusion: confirm that masked outer particles do not contribute to monopole moments in the tree and that pairwise P2P forces from outer sources are absent (beyond numerical tolerance);
\item MPI consistency: run identical initial conditions on 1 and on multiple ranks and verify that accelerations and early-time trajectories near the cutoff are numerically identical within roundoff;
\item restart repeatability: run to a snapshot, restart, and confirm the mask and potential behaviour are unchanged.
\end{itemize}

These checks are documented and included in the project's validation notes so the truncated-NFW backend can be accepted for experimental use only after passing the suite.

\subsection{Caveats and limitations}

The present implementation aims at a minimal, well-defined numerical experiment and therefore imposes several deliberate limitations:

\begin{itemize}
\item the cutoff is a hard step; no smooth transition shell is implemented (this can be added later if numerical artefacts at the boundary are observed),
\item only dark-matter particles are eligible for masking; gas, stars and sink/BH particles keep their full source mass, and their coupling to the external potential is unchanged,
\item using a fixed centre for the external potential while the particle distribution drifts can create a mismatch; the user must therefore choose the centre carefully (e.g. by using a shrinking-sphere centre estimate before launching a hybrid run),
\item energy is not strictly conserved across the particle/analytic split because masked particle mass is removed from the self-gravity source and replaced by an analytic potential; diagnostics are provided to quantify this effect,
\item for strict energy continuity diagnostics, a constant potential offset at $R_{\mathrm{cut}}$ may still be desirable in the external backend (forces are unaffected by this constant, but reported potential energies can shift).
\end{itemize}

Implementation details and exact file locations are documented in the project plan and inline in the code (see \texttt{software/swiftsim/src/potential/truncated\_nfw/potential.h} and related gravity hooks). The repository includes the configure option, the backend header, and the gravity-kernel adaptations required to run and validate the hybrid experiments described in Chapter~\ref{sec:framework_shells}.

